% Part: incompleteness
% Chapter: representability-in-q
% Section: undecidability

\documentclass[../../../include/open-logic-section]{subfiles}

\begin{document}

\olfileid{inc}{req}{und}
\olsection{عدم القابلية للقرار}

النظرية $\Th{T}$ \emph{غير قابلة للقرار} إذا امتنع وجود إجراء حسابي
ينتهي دائمًا في خطوات منتهية إلى جواب صحيح عن السؤال:
«أتثبت $\Th{T}$ الصيغة~$!A$؟»، وذلك لكل جملة~$!A$ في لغة
$\Th{T}$. فلو كانت $\Th{Q}$ قابلة للقرار لأمكن، بإجراء حسابي،
أخذ جملة~$!A$ من لغة الحساب والبتّ في صدق
$\Th{Q}\Proves !A$ أو عدمه، والعكس كذلك.
ويصير السؤال أدق إذا صغناه عن العلاقة~$\Prov[\Th{Q}](\OLId{latin-ordinary}{y})$،
التي تصدق على~$\OLId{latin-ordinary}{y}$ حين يكون $\OLId{latin-ordinary}{y}$ عدد غودل لجملة تثبتها~$\Th{Q}$:
أتكون هذه العلاقة عودية؟ سنثبت أن الجواب بالنفي.

\begin{thm}
يمتنع قرار $\Th{Q}$؛ وبعبارة أدق، العلاقة
\[
\Prov[\Th{Q}](\OLId{latin-ordinary}{y}) \defiff \fn{Sent}(\OLId{latin-ordinary}{y}) \land
\lexists[\OLId{latin-ordinary}{x}][\Prf[\Th{Q}](\OLId{latin-ordinary}{x}, \OLId{latin-ordinary}{y})]
\]
غير عودية.
\end{thm}

\begin{proof}
نفرض عودية هذه العلاقة، ونستخرج من الفرض حلًا لمسألة التوقف.
إذا أعطيت $\OLId{latin-ordinary}{e}$ و$\OLId{latin-ordinary}{n}$، فإن $\cfind{\OLId{latin-ordinary}{e}}(\OLId{latin-ordinary}{n})\fdefined$ تكافئ وجود~$\OLId{latin-ordinary}{s}$
تحقق $\OLId{latin-looped}{T}(\OLId{latin-ordinary}{e},\OLId{latin-ordinary}{n},\OLId{latin-ordinary}{s})$، حيث $\OLId{latin-looped}{T}$ محمول كليني المذكور في
\olref[cmp][rec][nft]{thm:kleene-nf}.
ولأن $\OLId{latin-looped}{T}$ عودي بدائي، تمثله في~$\Th{Q}$ صيغة $!B_\OLId{latin-looped}{T}$؛
فتكون $\Th{Q}\Proves !B_\OLId{latin-looped}{T}(\num{\OLId{latin-ordinary}{e}},\num{\OLId{latin-ordinary}{n}},\num{\OLId{latin-ordinary}{s}})$
إذا وفقط إذا صدقت $\OLId{latin-looped}{T}(\OLId{latin-ordinary}{e},\OLId{latin-ordinary}{n},\OLId{latin-ordinary}{s})$.
ومتى ثبتت $\Th{Q}\Proves !B_\OLId{latin-looped}{T}(\num{\OLId{latin-ordinary}{e}},\num{\OLId{latin-ordinary}{n}},\num{\OLId{latin-ordinary}{s}})$، ثبتت أيضًا
$\Th{Q}\Proves\lexists[\OLId{latin-ordinary}{y}][!B_\OLId{latin-looped}{T}(\num{\OLId{latin-ordinary}{e}},\num{\OLId{latin-ordinary}{n}},\OLId{latin-ordinary}{y})]$.
أما إذا لم يوجد مثل هذا $\OLId{latin-ordinary}{s}$، فإن
$\Th{Q}\Proves\lnot !B_\OLId{latin-looped}{T}(\num{\OLId{latin-ordinary}{e}},\num{\OLId{latin-ordinary}{n}},\num{\OLId{latin-ordinary}{s}})$ لكل~$\OLId{latin-ordinary}{s}$.
والنظرية $\Th{Q}$ متسقة اتساقًا من نوع $\omega$؛
أي إنه متى كانت $\Th{Q}\Proves\lnot !A(\num{\OLId{latin-ordinary}{n}})$
لكل~$\OLId{latin-ordinary}{n}\in\Nat$، لزم $\Th{Q}\Proves/\lexists[\OLId{latin-ordinary}{y}][!A(\OLId{latin-ordinary}{y})]$.
ودليل ذلك صدق بديهيات $\Th{Q}$ في النموذج القياسي~$\Struct{N}$.
فينتج $\Th{Q}\Proves/\lexists[\OLId{latin-ordinary}{y}][!B_\OLId{latin-looped}{T}(\num{\OLId{latin-ordinary}{e}},\num{\OLId{latin-ordinary}{n}},\OLId{latin-ordinary}{y})]$.
فالمحصّل أن $\Th{Q}\Proves\lexists[\OLId{latin-ordinary}{y}][!B_\OLId{latin-looped}{T}(\num{\OLId{latin-ordinary}{e}},\num{\OLId{latin-ordinary}{n}},\OLId{latin-ordinary}{y})]$
تكافئ وجود $\OLId{latin-ordinary}{s}$ تحقق $\OLId{latin-looped}{T}(\OLId{latin-ordinary}{e},\OLId{latin-ordinary}{n},\OLId{latin-ordinary}{s})$، أي تكافئ
$\cfind{\OLId{latin-ordinary}{e}}(\OLId{latin-ordinary}{n})\fdefined$.
ومن $\OLId{latin-ordinary}{e}$ و$\OLId{latin-ordinary}{n}$ يمكن حساب
$\Gn{\lexists[\OLId{latin-ordinary}{y}][!B_\OLId{latin-looped}{T}(\num{\OLId{latin-ordinary}{e}},\num{\OLId{latin-ordinary}{n}},\OLId{latin-ordinary}{y})]}$؛
فلنسمّ الدالة العودية البدائية التي تقوم بذلك $\OLId{latin-ordinary}{g}(\OLId{latin-ordinary}{e},\OLId{latin-ordinary}{n})$.
ونعرّف:
\[
\OLId{latin-ordinary}{h}(\OLId{latin-ordinary}{e}, \OLId{latin-ordinary}{n}) =
\begin{cases}
\OLMathNumeral{1} & \text{إذا كانت $\Prov[\Th{Q}](g(e, n))$ صادقة}\\
\OLMathNumeral{0} & \text{في غير ذلك}.
\end{cases}
\]
فتكون $\OLId{latin-ordinary}{h}$ عودية على فرض عودية $\Prov[\Th{Q}]$.
لكن $\OLId{latin-ordinary}{h}$ غير عودية بمقتضى \olref[cmp][rec][hlt]{thm:halting-problem}؛
فبطل الفرض، وامتنع أن تكون $\Prov[\Th{Q}]$ عودية.
\end{proof}

\begin{cor}
يمتنع وجود إجراء قرار لمنطق الرتبة الأولى.
\end{cor}

\begin{proof}
لو أمكن قرار منطق الرتبة الأولى لأمكن قرار قابلية الإثبات في
$\Th{Q}$؛ إذ إن $\Th{Q}\Proves !A$ تكافئ
$\Proves !T\lif !A$، حيث $!T$ اقتران بديهيات~$\Th{Q}$ المنتهية.
وقد ثبت امتناع قرارها.
\end{proof}

\end{document}
